%
% Portada.
%

% Nota: Ser�a m�s c�modo emplear el comando \maketitle que genera una portada de forma autom�tica, pero 
% no incluye toda la informaci�n que es necesario incluir en la memoria de un proyecto de fin de carrera
% de la Facultad de Inform�tica de A Coru�a.
%

\begin{titlepage}

	\begin{center}

		% Logotipo de la universidad.
		\includegraphics[width=6cm]{eps/00.eps}
		\vspace{2cm}

		% Nombre de la facultad, de la universidad y del departamento en que se realiza el PFC.
		{\Large{\textbf{Faculty of Computer Science}}}
		\\
		{\Large{\textbf{University of A Coru�a}}}
		\vspace{1cm}

		% Indicamos el nombre de la titulaci�n oficial que hemos cursado con tanto esfuerzo.
		{\large FINAL YEAR PROJECT\\COMPUTER SCIENCE DEGREE}
		\vspace{1cm}

		% T�tulo
		\textbf{\LARGE Real-time management tool for massive 3D point clouds}
		\vspace{7cm}
	\end{center}

	\begin{flushright}
		\begin{tabular}{ll}
			% Nombre del alumno.
			\large{\textbf{Student:}}	&
			\large{Luis Omar Alvarez Mures} \\

			% Nombre del director/tutor del proyecto.
			\large{\textbf{Advisors:}}	&
			\large{Emilio Jos� Padr�n Gonz�lez} \\
			 & \large{Alberto Jaspe Villanueva} \\

			% Fecha.
			\large{\textbf{Date:}}	&
			\large{\today} \\
		\end{tabular}
	\end{flushright}

\end{titlepage}
